% 附录 B LVL Utility
\biappendix{LVL 实用程序}{LVL Utility}
\label{app:b}

本附录介绍 IC Validator 工具提供的版图对版图（Layout-Versus-Layout，LVL）实用程序以及 QuickLVL 实用程序。

LVL 实用程序包括以下各节：
\begin{itemize}
  \item 说明（Description）
  \item 命令行选项（Command-Line Options）
  \item 使用 OpenAccess 库格式（Using the OpenAccess Library Format）
  \item 层文件（Layer Files）
  \item 比较层子集（Comparing a Subset of Layers）
  \item 文本处理（Handling of Text）
  \item LVL 输出（LVL Output）
  \item 使用模型（Use Models）
\end{itemize}

QuickLVL 实用程序包括以下各节：
\begin{itemize}
  \item QuickLVL 实用程序（QuickLVL Utility）
  \item 单元级 Quick LVL（Cell-Level Quick LVL）
  \item LVL 与 QuickLVL 的差异（Differences Between LVL and QuickLVL）
  \item 聚焦差异层与区域的自动增量 DRC（Auto-Incremental DRC Focused on Differing Layers and Regions）
\end{itemize}

% ============================================================
\section{说明}
\subsection*{Description}

LVL 实用程序可用于：
\begin{itemize}
  \item 从单个库生成 assign 文件。
  \item 使用基于 XOR 的比较或基于 NOT 的比较，对设计的两个版图进行比较。发现的错误可通过 \cmd{cell1\_vs\_cell2.vue} 文件在 VUE 中查看。更多信息见 \textit{IC Validator VUE User Guide}。
\end{itemize}

\begin{noteBox}
LVL 实用程序支持除 NULL 字符外的任意文本字符。
\end{noteBox}

对于基于 NOT 的比较：
\begin{itemize}
  \item 基线层与检查库层做 NOT 运算；
  \item 检查库层与基线层做 NOT 运算；
  \item 各 NOT 运算的错误分别显示在 \cmd{cell1\_vs\_cell2.LVL\_ERRORS} 文件中。
\end{itemize}

LVL 实用程序支持下列版图输入格式：
\begin{itemize}
  \item GDSII
  \item Milkyway
  \item OASIS
  \item OpenAccess
  \item NDM
\end{itemize}

\begin{noteBox}
LVL 实用程序支持除 NULL 字符外的任意文本字符。
\end{noteBox}

% ============================================================
\section{命令行选项}
\subsection*{Command-Line Options}

LVL 实用程序的命令行语法为：

\begin{lstlisting}
icv_lvl library1 library2 -c cell | -c1 cell1 -c2 cell2 [options]
\end{lstlisting}

表~\ref{tab:app-b-lvl-opts} 说明基本命令行选项以及可用于 LVL 分布式处理的命令行选项。

\begin{longtable}{@{}>{\ttfamily\raggedright\arraybackslash}p{0.32\textwidth} p{0.62\textwidth}@{}}
\caption{LVL 实用程序命令行选项}\\
\label{tab:app-b-lvl-opts}\\
\toprule
\textbf{\rmfamily 语法} & \textbf{说明} \\
\midrule
\endfirsthead
\toprule
\textbf{\rmfamily 语法} & \textbf{说明} \\
\midrule
\endhead
\bottomrule
\endfoot

-64 &
以支持 64 位坐标的方式运行该实用程序。 \\
\midrule

-af \textit{filename} &
存放库中 assign 层的文件名。该文件由 LVL 实用程序生成，语法与 IC Validator runset 的 ASSIGN 节相同。若希望文件不在当前目录，请包含路径。$^{a}$ \\
\midrule

-c \textit{cell} &
两个库 \cmd{library1} 与 \cmd{library2} 的单元名。使用 \cmd{-af} 生成 assign 文件时也需要单元名。$^{a}$ \\
\midrule

-c1 \textit{cell1} &
第一个库 \cmd{library1} 的单元名。$^{a}$ \\
\midrule

-c2 \textit{cell2} &
第二个库 \cmd{library2} 的单元名。 \\
\midrule

-cell\_map\_file &
支持将第二个输入库中的单元名映射到第一个输入库中的单元名。 \\
\midrule

-elpc \textit{num}\newline
-error\_limit\_per\_check \textit{num} &
设置每类检查的错误上限。默认值为 100。
\textit{注意：}将该选项设为 \cmd{ERROR\_LIMIT\_MAX} 时，将存储直至工具定义的最大上限的全部错误。也可给出具体数字，将上限设为高于工具定义的最大值；但对大量错误，提高上限可能导致性能下降与磁盘占用增加，因此不推荐。设为 0 则抑制存储全部错误。 \\
\midrule

-flat &
将全部文本与几何展平到顶层单元，并在顶层单元中执行所有层比较。 \\
\midrule

-flatten &
将全部输出展平到顶层单元。 \\
\midrule

-h &
显示用法信息。 \\
\midrule

-host\_add \textit{hostname}\newline
\textit{| number\_of\_cpus |}\newline
\textit{hostname:number\_of\_cpus | LSF | SGE | NB} &
向已在运行的 \cmd{icv\_lvl} 作业添加更多主机（含远程主机）。用法与 \cmd{-host\_init} 相同。详见「指定主机名与 CPU 数量」。 \\
\midrule

-host\_init &
在单主机或多主机上启动 IC Validator 作业。详见「指定主机名与 CPU 数量」。登录远程主机所用远程登录协议的优先级见 \cmd{-host\_login}。更多信息见第~\ref{chap:basics}~章中 IC Validator 的 \cmd{-host\_init} 命令行选项。 \\
\midrule

-host\_init LSF &
通过 LSF 批处理排队系统获取主机，读取 \cmd{LSB\_HOSTS} 或 \cmd{LSB\_MCPU\_HOSTS} 环境变量。远程登录协议优先级见 \cmd{-host\_login}。更多信息见第~\ref{chap:basics}~章中 \cmd{-host\_init LSF}。 \\
\midrule

-host\_init SGE &
通过 SGE 批处理排队系统获取主机，读取 \cmd{PE\_HOSTFILE} 环境变量。远程登录协议优先级见 \cmd{-host\_login}。更多信息见第~\ref{chap:basics}~章中 \cmd{-host\_init SGE}。 \\
\midrule

-host\_disk &
为 \cmd{-host\_init} 或 \cmd{-host\_add} 指定的全部主机指定本地组目录。更多信息见第~\ref{chap:basics}~章中 \cmd{-host\_disk}。 \\
\midrule

-host\_login &
使 IC Validator 能登录远程主机以启动分布式作业进程。请确保工具能找到提供给 \cmd{-host\_login} 的二进制/脚本。默认为 \cmd{rsh}。更多信息见第~\ref{chap:basics}~章中 \cmd{-host\_login}。 \\
\midrule

-host\_memory \textit{Memory\_Value} M|G|T &
通过重调度与线程化，尝试使作业所用主机维持用户指定的目标内存。不能用该开关降低单个命令的内存。更多信息见第~\ref{chap:basics}~章中 \cmd{-host\_memory}。 \\
\midrule

-ignore\_text\_case &
指示实用程序在比较文本时忽略大小写。默认行为是文本比较区分大小写。 \\
\midrule

-al &
与 \cmd{-layer\_file} 参数一起使用。向 LVL 运行添加任何缺失的层。因需扫描输入文件以判断层文件是否缺失层，会有性能影响。仅当被检查文件为 GDSII 或 OASIS 时可用。 \\
\midrule

-lf \textit{filename}\newline
-layer\_file \textit{filename} &
两个库共用的层文件，用于层映射、层合并与层命名。必须包含有效的 PXL runset ASSIGN 节。若库不在当前目录，请包含路径。更多信息见「层文件」。$^{b}$ \\
\midrule

-lf1 \textit{filename}\newline
-layer\_file1 \textit{filename} &
第一个库的层文件。始终与 \cmd{-lf2} 一起使用。若库不在当前目录，请包含路径。更多信息见「层文件」。$^{b}$ \\
\midrule

-lf2 \textit{filename}\newline
-layer\_file2 \textit{filename} &
第二个库的层文件。始终与 \cmd{-lf1} 一起使用。若库不在当前目录，请包含路径。更多信息见「层文件」。$^{b}$ \\
\midrule

-oa\_color\_mapping \textit{option} &
指定如何使用锁定/未锁定锚定位（anchor bit）进行 OpenAccess 颜色层映射。\textit{option} 可选：
\begin{itemize}[nosep,leftmargin=1.2em]
  \item \cmd{ignore}：使用未指定 \cmd{locked}/\cmd{unlocked} 关键字的颜色映射，忽略已指定者。
  \item \cmd{strict}：使用已指定 \cmd{locked}/\cmd{unlocked} 的颜色映射，忽略未指定者。
  \item \cmd{full}：将读入时带颜色的全部形状设为 locked。映射与 \cmd{compatibility} 相同。
  \item \cmd{compatibility}：遵循行业标准。忽略颜色规范之后出现 \cmd{colorAnnotate} 关键字的颜色映射。若无规则匹配彩色形状，则使用无颜色的映射。
  \item \cmd{lockOnly}：对可着色层，若技术数据库中该层的 \cmd{-fullColor} 为 true（默认），则仅映射已着色且锁定的形状；未着色或未锁定的形状丢弃。若某层 \cmd{-fullColor} 为 false，则对该层遵守层映射。
\end{itemize} \\
\midrule

-oa\_color\_mapping1 ignore | strict | full | compatibility | lockOnly &
针对 OpenAccess 库1 的颜色层映射形式，指定 locked/unlocked 锚定位用法。 \\
\midrule

-oa\_color\_mapping2 ignore | strict | full | compatibility | lockOnly &
针对 OpenAccess 库2 的颜色层映射形式，指定 locked/unlocked 锚定位用法。 \\
\midrule

-of \textit{options\_file} &
可选文件，包含 IC Validator runset 的 OPTIONS 节。可含任意 OPTIONS 节函数，包括 \cmd{gds\_options()}、\cmd{milkyway\_options()} 与 \cmd{oasis\_options()}。 \\
\midrule

-rsi &
报告下列输入库信息。仅支持 GDSII 与 OASIS 输入格式。报告内容包括：GDSII 最后修改与最后访问时间；用 \cmd{md5sum -b file} 计算的 md5sum；输入库格式；输入库路径；用于库输入的命令。枚举须为逗号分隔列表，且整个列表放在引号内。见 \cmd{run\_options()} 函数的 \cmd{report\_streamfile\_information} 参数。
\textit{注意：}\cmd{-rsi} 命令行选项会覆盖 \cmd{run\_options()} 中的 runset 设置。 \\
\midrule

\textit{library1} &
基线库名。若不在当前目录请包含路径。当用 \cmd{-af} 只指定一个库时，LVL 实用程序生成 assign 文件。不能在同一次 LVL 运行中既生成 assign 文件又比较两个版图。各格式输入示例：
\begin{itemize}[nosep,leftmargin=1.2em]
  \item GDSII：\cmd{/path/to/file.gds}
  \item OASIS：\cmd{/path/to/file.oas}
  \item Milkyway：\cmd{/library/path/LIBRARY\_NAME}
  \item OpenAccess：\cmd{/path/to/lib.defs/LIBRARY\_NAME}
  \item NDM：\cmd{/library/path/LIBRARY\_NAME}
\end{itemize} \\
\midrule

\textit{library2} &
检查库名。若不在当前目录请包含路径。使用 \cmd{-af} 时不要指定该文件。不能在同一次运行中既生成 assign 文件又比较两个版图。输入格式示例同 \textit{library1}。 \\
\midrule

-mf \textit{filename}\newline
-map\_file \textit{filename} &
应用于两个库的层映射文件。更多信息见「映射层」。 \\
\midrule

-mf1 \textit{filename}\newline
-map\_file1 \textit{filename} &
应用于第一个库的层映射文件。更多信息见「映射层」。 \\
\midrule

-mf2 \textit{filename}\newline
-map\_file2 \textit{filename} &
应用于第二个库的层映射文件。更多信息见「映射层」。 \\
\midrule

-ndl \textit{label}\newline
-ndm\_design\_label &
两个库 \cmd{library1} 与 \cmd{library2} 的设计标签。$^{a}$ \\
\midrule

-ndl1 \textit{label}\newline
-ndm\_design\_label1 &
第一个库 \cmd{library1} 的设计标签。$^{a}$ \\
\midrule

-ndl2 \textit{label}\newline
-ndm\_design\_label2 &
第二个库 \cmd{library2} 的设计标签。 \\
\midrule

-not &
指示实用程序做两次基于 NOT 的比较，而非基于 XOR 的比较。 \\
\midrule

-oa\_dm4 &
使用与 DM4 插件兼容的 OpenAccess 共享库。 \\
\midrule

-outside &
指示实用程序做两次基于 OUTSIDE 的比较，而非基于 XOR 的比较。执行 \cmd{not\_interacting()} 时，任何接触的数据不纳入该运算。可与 \cmd{-report\_touch} 一起使用，使接触被视为不交互。 \\
\midrule

-quick &
执行单元级 Quick LVL。更多信息见「QuickLVL 实用程序」。
\textit{注意：}Quick LVL（\cmd{-quick}）仅在单主机上工作。若从批处理排队系统获取了多主机，Quick LVL 将报错退出。 \\
\midrule

-resolution &
允许在 LVL 实用程序内指定 IC Validator 运行的内部分辨率。 \\
\midrule

-rt ALL | POINT\newline
-report\_touch ALL | POINT &
指定如何处理点接触。与 \cmd{-outside} 一起使用。
\begin{itemize}[nosep,leftmargin=1.2em]
  \item \cmd{ALL}：将多边形线接触与边点接触视为 outside。
  \item \cmd{POINT}：将点接触视为 outside。
\end{itemize} \\
\midrule

-sel "\textit{layer1[-layer1max][,dtype1[-dtype1max]] [...]}"\newline
-select\_edge\_layers &
指定比较所用的边层与 datatype。可使用范围。未指定 datatype 时默认为全部 datatype。$^{c}$ \\
\midrule

-small\_cell &
强制小单元模式，使 LVL 比较部分在小设计上尽可能快。若工具判定输入设计过大，为避免在大设计上变慢，将以错误退出。 \\
\midrule

-sn \textit{layer\_name1} \ldots\newline
-select\_names \textit{layer\_name1} \ldots &
按名称指定实用程序所运行的层。 \\
\midrule

-spl "\textit{layer1[-layer1max][,dtype1[-dtype1max]] [...]}"\newline
-select\_polygon\_layers &
指定比较所用的多边形层与 datatype。可使用范围。未指定 datatype 时默认为全部 datatype。要选择文本层进行比较，应使用 \cmd{-stl}。$^{c}$ \\
\midrule

-split\_output &
与 \cmd{-not} 一起使用时，将 not12 结果写入 \cmd{Cell1\_vs\_Cell2\_not12}，将 not21 结果写入 \cmd{Cell1\_vs\_Cell2\_not21}。not12 对应对基线库（命令行第一个库）各层与检查库对应层执行 (\textit{BaselineLayer} not \textit{Checklayer})；not21 对应 (\textit{CheckLayer} not \textit{BaselineLayer})。 \\
\midrule

-stl "\textit{layer1[-layer1max][,dtype1[-dtype1max]] [...]}"\newline
-select\_text\_layers &
指定比较所用的文本层与 datatype。可使用范围。未指定 datatype 时默认为全部 datatype。要选择多边形层，应使用 \cmd{-spl}。$^{c}$$^{d}$ \\
\midrule

-sub\_cell &
指示对较大 OASIS 版图内的小下级单元运行，避免读取整个版图。 \\
\midrule

-suppress\_empty\_layer\_errors &
抑制报告空层错误多边形。 \\
\midrule

-text &
指示实用程序比较两个版图之间的文本。更多信息见「文本处理」。未使用该选项时，仅比较几何，不比较文本。 \\
\midrule

-top\_cell\_only &
仅比较各库顶层单元中的数据，丢弃全部下层单元数据。 \\
\midrule

-uel "\textit{layer1[-layer1max][,dtype1[-dtype1max]] [...]}"\newline
-unselect\_edge\_layers &
指定从比较中排除的边层与 datatype。可使用范围。未指定 datatype 时默认为全部 datatype。$^{c}$ \\
\midrule

-un \textit{layer\_name1} \ldots\newline
-unselect\_names \textit{layer\_name1} \ldots &
按名称指定运行时排除的层。 \\
\midrule

-upl "\textit{layer1[-layer1max][,dtype1[-dtype1max]] [...]}"\newline
-unselect\_polygon\_layers &
指定从比较中排除的多边形层与 datatype。可使用范围。未指定 datatype 时默认为全部 datatype。$^{c}$ \\
\midrule

-utl "\textit{layer1[-layer1max][,dtype1[-dtype1max]] [...]}"\newline
-unselect\_text\_layers &
指定从比较中排除的文本层与 datatype。可使用范围。未指定 datatype 时默认为全部 datatype。$^{c}$$^{d}$ \\
\midrule

-transform [RX,][0|90|180|270|DX,DY] &
指示按命令行描述变换检查版图的顶层单元。RX：绕 X 轴反射；0|90|180|270：逆时针旋转角；dx dy：X/Y 方向平移（微米）。示例：\cmd{-transform 90,0,0}、\cmd{-transform RX,90,1,-1.2}、\cmd{-transform RX}。更多信息见「QuickLVL 实用程序」。 \\
\midrule

\textit{baseline\_layout} &
基线库名（QuickLVL）。若不在当前目录请包含路径。示例：OASIS \cmd{/path/to/file.oas}；GDSII \cmd{/path/to/file.gds}。更多信息见「QuickLVL 实用程序」。 \\
\midrule

\textit{check\_layout} &
检查库名（QuickLVL）。若不在当前目录请包含路径。格式示例同基线。更多信息见「QuickLVL 实用程序」。 \\
\midrule

-V &
打印实用程序版本。 \\
\end{longtable}

{\footnotesize
\begin{itemize}[leftmargin=1.5em,nosep]
  \item[$^{a}$] 若两库单元名相同，用 \cmd{-c} 提供单元名；若不同，必须用 \cmd{-c1} 与 \cmd{-c2} 分别提供。
  \item[$^{b}$] 层文件可以是完整 runset；LVL 实用程序忽略 ASSIGN 节以外的全部内容。
  \item[$^{c}$] 若比较两个 OpenAccess 库，可用层名与 purpose 名代替层号与 datatype 号。
  \item[$^{d}$] 仅当位置与文本字符串均匹配时文本才匹配。文本差异在 VUE 中显示为 X。文本比较仅为单元级操作。层次不匹配时，先比较层次与版图，再比较文本；若层次与版图不匹配，则无法准确比较文本。一般而言，若两个单元的名称、尺寸与放置不同，其文本不会被比较。
\end{itemize}
}

% ============================================================
\section{使用 OpenAccess 库格式}
\subsection*{Using the OpenAccess Library Format}

比较 OpenAccess 格式版图时，下列命令行选项很有用：
\begin{itemize}
  \item \cmd{-oalm} / \cmd{-oa\_layer\_map} \textit{filename}
  \item \cmd{-oalm1} / \cmd{-oa\_layer\_map1} \textit{filename}
  \item \cmd{-oalm2} / \cmd{-oa\_layer\_map2} \textit{filename}
  \item \cmd{-oacm} / \cmd{-oa\_cell\_map} \textit{cellmap}
  \item \cmd{-oacm1} / \cmd{-oa\_cell\_map1} \textit{cellmap1}
  \item \cmd{-oacm2} / \cmd{-oa\_cell\_map2} \textit{cellmap2}
  \item \cmd{-oa\_color\_mapping ignore | strict | full | compatibility}
  \item \cmd{-oa\_color\_mapping1 ignore | strict | full | compatibility}
  \item \cmd{-oa\_color\_mapping2 ignore | strict | full | compatibility}
  \item \cmd{-oav} / \cmd{-oa\_view} \textit{viewname}
  \item \cmd{-oav1} / \cmd{-oa\_view1} \textit{viewname}
  \item \cmd{-oav2} / \cmd{-oa\_view2} \textit{viewname}
\end{itemize}

若比较两个 OpenAccess 库，无需指定层映射文件；LVL 实用程序按 layer/purpose 对匹配并运行比较。但若将 OpenAccess 库与 GDSII、OASIS、Milkyway 或 NDM 库比较，则需要层映射文件，以将 OpenAccess 的 layer/purpose 映射到 GDSII/Milkyway/OASIS/NDM 的 layer/datatype。例如：

\begin{lstlisting}
% icv_lvl lib.defs/OA_LIB gds_lib.gds -c top -oalm layer_map.txt
\end{lstlisting}

在本例中：
\begin{itemize}
  \item 先列出的 OpenAccess 库（\cmd{OA\_LIB}）与 GDSII 库（\cmd{gds\_lib.gds}）比较。
  \item \cmd{lib.defs} 文件列出设计所需的全部库（含 OpenAccess 库）。文件名必须为 \cmd{lib.defs}。若不在当前目录，请给出路径并接文件名，例如 \cmd{datapath1/testcase1/lib.defs/OA\_LIB}。
\end{itemize}

\cmd{lib.defs} 文件内容示例：
\begin{lstlisting}
DEFINE analogLib ./PDK/lib/analogLib
ASSIGN analogLib libMode shared
DEFINE basic ./PDK/lib/basic
ASSIGN basic libMode shared
DEFINE device_lib ./PDK/lib/device_lib
ASSIGN device_lib libMode shared
DEFINE OA_LIB ./OA_LIB
\end{lstlisting}

\subsection{OpenAccess 层映射文件}
\subsubsection*{OpenAccess Layer Mapping File}

若比较两个 OpenAccess 数据库，无需提供层映射文件；LVL 按 layer/purpose 对匹配并比较。

但若将 OpenAccess 数据库与 GDSII、OASIS、Milkyway 或 NDM 库比较，则需提供层映射文件，将 OpenAccess 的 layer/purpose 映射到目标文件的 layer/datatype。

若两个 OpenAccess 库使用同一层映射文件，指定：
\begin{lstlisting}
-oalm filename
\end{lstlisting}

若各输入 OA 数据库有独立层映射文件，分别指定：
\begin{lstlisting}
-oalm1 filename1
-oalm2 filename2
\end{lstlisting}

更多信息见 \textit{IC Validator Reference Manual} 中 \cmd{openaccess\_options()} 函数的 \cmd{layer\_mapping\_file} 参数。

OpenAccess 层映射格式示例：
\begin{lstlisting}
#OA Layer Name  OA Purpose Name  Stream layer  Stream Datatype
Layer1          drawing          1             0
Layer2          drawing          2             1
Layer3          drawing          3             2
Layer4          drawing          4             0
Layer4          boundary         4             1
\end{lstlisting}

\subsection{单元映射}
\subsubsection*{Cell Mapping}

将 OpenAccess 数据库与其他格式比较时，可能需要提供单元映射以防止展平。OpenAccess 库用 library/cell/view 三元组标识唯一单元。IC Validator 与 LVL 必须为每个唯一三元组分配唯一单元名。LVL 会展平单元名不匹配的层次，直到两层次匹配。使用单元映射文件可保证被比较两库的单元命名一致，从而避免展平。

更多信息见 \textit{IC Validator Reference Manual} 中 \cmd{openaccess\_options()} 的 \cmd{cell\_mapping\_file} 参数。

若两个 OpenAccess 库使用同一单元映射，指定：
\begin{lstlisting}
-oacm cellmap
\end{lstlisting}

若各输入数据库有独立映射，分别指定：
\begin{lstlisting}
-oacm1 cellmap1
-oacm2 cellmap2
\end{lstlisting}

单元映射文件示例：
\begin{lstlisting}
#Library Name  Cell Name  View Name  Stream Name
des_lib        nand2      layout     nand2_1
sdes_lib       nand2      layout     nand2_2
des_lib        cell_nu    layout     cell_nu
sdes_lib       cell_nu    layout     cell_new
\end{lstlisting}

\subsection{OpenAccess 基线库的错误输出格式}
\subsubsection*{Error Output Format for OpenAccess Baseline Library}

错误输出写入当前工作目录下名为 \cmd{cell1\_vs\_cell2} 的 OpenAccess 库。

输出库中的错误数据尽可能使用基线库的层名。

未给出层文件时，层名一般格式为：
\begin{lstlisting}
LlayerPpurpose
\end{lstlisting}

输出 purpose 名为 \cmd{Xor}、\cmd{BaselineNotCheck} 或 \cmd{CheckNotBaseline}，取决于所选 LVL 方法。

输出示例：
\begin{lstlisting}
Checking LxxxPyyy_Baseline XOR LaaaPbbb_Check, output to (LxxxPyyy; Xor)
Checking LxxxPyyy_Baseline NOT LaaaPbbb_Check, output to (LxxxPyyy; BaselineNotCheck)
Checking LaaaPbbb_Check NOT LxxxPyyy_Baseline, output to (LaaaPbbb; CheckNotBaseline)
\end{lstlisting}

\subsection{设置单元视图名}
\subsubsection*{Setting the Cell View Name}

在命令行指定 OpenAccess 库时，可能需要将顶层单元的视图名设为默认 \cmd{layout} 以外的值。

若两个 OpenAccess 库使用同一视图名，指定：
\begin{lstlisting}
-oav viewname
\end{lstlisting}

若各输入数据库顶层单元视图名不同，分别指定：
\begin{lstlisting}
-oav1 viewname1
-oav2 viewname2
\end{lstlisting}

% ============================================================
\section{层文件}
\subsection*{Layer Files}

默认比较版图数据库的全部层。通过命令行选项，可对层子集运行实用程序；也可指定映射文件以替代 Milkyway 与 NDM 库内部包含的映射文件。

层先经映射，再经命令行选项选择或过滤。

例如，若同时使用 \cmd{-lf layerfile} 与 \cmd{-mf mapfile}，则先用 \cmd{mapfile} 映射层号与 datatype，再解释 \cmd{layerfile}。

\subsection{选择层子集}
\subsubsection*{Selecting a Subset of Layers}

使用 \cmd{-lf}，或同时使用 \cmd{-lf1} 与 \cmd{-lf2} 指定层文件时，仅比较层文件中出现的层；未出现在任何 assign 语句中的层号不做比较。

层文件限制如下：
\begin{itemize}
  \item 当同一 layer/datatype 对多次使用时，冗余 assign 被忽略。下例中第三个赋值 \cmd{L3D0} 被忽略：
\begin{lstlisting}[style=pxl]
L1D0 = assign({{1,0}});
L2D0 = assign({{2,0}});
L3D0 = assign({{1,0}});
\end{lstlisting}
  \item 若层文件含同一层名的多个定义，LVL 报错停止。下例中层 A 定义两次，须删除其一：
\begin{lstlisting}[style=pxl]
A = assign({{1,0}});
B = assign({{2,0}});
A = assign({{10,0}});
\end{lstlisting}
\end{itemize}

\subsection{映射层}
\subsubsection*{Mapping Layers}

映射文件可作为 Milkyway 与 NDM 库的内部部分。默认使用内部映射文件。用 \cmd{-mf}、\cmd{-mf1}、\cmd{-mf2} 可通过外部文件映射层；命令行指定的映射文件将替代内部映射文件。

\begin{noteBox}
要对 OpenAccess 库用 \cmd{-mf}/\cmd{-mf1}/\cmd{-mf2} 映射层，还必须用 \cmd{-oalm}/\cmd{-oalm1}/\cmd{-oalm2} 将 LayerName/Purpose 字符串对转换为 LayerNumber/Datatype 数字对。
\end{noteBox}

映射文件可为下列三种格式之二：Milkyway 格式、NDM 格式或 IC Validator 映射文件格式。Milkyway 与 NDM 格式见 SolvNetPlus 上的 \textit{Milkyway Database Application Note} 与 \textit{Library Data Preparation for IC Compiler User Guide}。IC Validator 映射文件格式见「层映射」。

\cmd{-lf}、\cmd{-lf1}、\cmd{-lf2}、\cmd{-spl}、\cmd{-upl}、\cmd{-sel}、\cmd{-uel}、\cmd{-stl}、\cmd{-utl}、\cmd{-sn} 与 \cmd{-un} 在层映射之后操作。

若不希望使用内部映射文件，可用映射文件命令行选项关闭：\cmd{-mf1 ""} 关闭第一个库；\cmd{-mf2 ""} 关闭第二个库；\cmd{-mf ""} 关闭两个库。

例如，两个输入库：
\begin{itemize}
  \item \cmd{one.gds}：数据在层 1、datatype 0；
  \item \cmd{two.gds}：数据在层 2、datatype 0，与 \cmd{one.gds} 层 1 数据相同。
\end{itemize}

\cmd{mf2} 映射文件含：
\begin{lstlisting}
D 2:0 1 0
\end{lstlisting}

LVL 命令为：
\begin{lstlisting}
icv_lvl one.gds two.gds -mf2 mf2 -c top
\end{lstlisting}

\cmd{-mf2} 将 \cmd{two.gds} 中层 2 数据映射到层 1，使两库数据落在同一层，LVL 比较干净。

也可用层文件达到同样效果：
\begin{itemize}
  \item \cmd{file1} 含 \cmd{Metal1 = assign(\{\{1, 0\}\});}
  \item \cmd{file2} 含 \cmd{Metal1 = assign(\{\{2, 0\}\});}
\end{itemize}

\begin{lstlisting}
icv_lvl one.gds two.gds -lf1 file1 -lf2 file2 -c top
\end{lstlisting}

若已有可用于 LVL 的 IC Validator 层映射文件，第一种（映射文件）方法更简便。

% ============================================================
\section{比较层子集}
\subsection*{Comparing a Subset of Layers}

可用过滤命令行选项比较多边形与边层的子集（可有或无层文件）。规则：
\begin{itemize}
  \item 不能将 select 选项与对应的 unselect（排除）选项同时使用；
  \item 不能将按名过滤与按号过滤同时使用。
\end{itemize}

过滤选项见表~\ref{tab:app-b-filter}。

\begin{table}[htbp]
\centering
\caption{过滤选项}
\label{tab:app-b-filter}
\begin{tabular}{@{}lll@{}}
\toprule
数据过滤类别 & 按名称 & 按编号 \\
\midrule
Select & \cmd{-sn} & \cmd{-spl}、\cmd{-sel}、\cmd{-stl} \\
Unselect & \cmd{-un} & \cmd{-upl}、\cmd{-uel}、\cmd{-utl} \\
\bottomrule
\end{tabular}
\end{table}

\begin{noteBox}
同一次运行中不能同时使用 \cmd{-sn} 与 \cmd{-un}；不能同时使用 \cmd{-spl} 与 \cmd{-upl}；不能同时使用 \cmd{-sel} 与 \cmd{-uel}；不能同时使用 \cmd{-stl} 与 \cmd{-utl}。
\end{noteBox}

\subsection{按名称比较}
\subsubsection*{Comparing by Name}

在用 \cmd{-lf} 或同时用 \cmd{-lf1}/\cmd{-lf2} 提供层文件时，可用命令行选项限制所比较的多边形与边层。此为按名过滤。

按名限制层的选项为：
\begin{itemize}
  \item \cmd{-sn}/\cmd{-select\_names} \textit{layer\_name1} \ldots
  \item \cmd{-un}/\cmd{-unselect\_names} \textit{layer\_name1} \ldots
\end{itemize}

示例：在 POLY、NWELL、METAL1 上运行：
\begin{lstlisting}
-sn POLY NWELL METAL1
\end{lstlisting}

在除 METAL2 外的全部层上运行：
\begin{lstlisting}
-un METAL2
\end{lstlisting}

\subsection{按编号比较}
\subsubsection*{Comparing by Number}

未提供层文件时，可用命令行选项限制所比较的层。此为按号过滤。多边形、边与文本层分别过滤。

按号限制层的选项为：
\begin{itemize}
  \item \cmd{-sel}/\cmd{-select\_edge\_layers} \texttt{"layer1[-layer1max][,dtype1[-dtype1max]] [...]"}
  \item \cmd{-uel}/\cmd{-unselect\_edge\_layers} \texttt{"\ldots"}
  \item \cmd{-spl}/\cmd{-select\_polygon\_layers} \texttt{"\ldots"}
  \item \cmd{-upl}/\cmd{-unselect\_polygon\_layers} \texttt{"\ldots"}
  \item \cmd{-stl}/\cmd{-select\_text\_layers} \texttt{"\ldots"}
  \item \cmd{-utl}/\cmd{-unselect\_text\_layers} \texttt{"\ldots"}
\end{itemize}

下例在下列多边形层上运行：层 1--4 全部 datatype；层 5 datatype 0--7；层 6 全部 datatype；层 7--8 datatype 0--7：
\begin{lstlisting}
-spl 1-4  5,0-7  6  7-8,0-7
\end{lstlisting}

下例在除多边形层 1、2、3 的 datatype 5、6、7 外的全部多边形层与 datatype 上运行：
\begin{lstlisting}
-upl 1-3,5-7
\end{lstlisting}

边层示例：层 1--5 datatype 0--7；层 6--8 datatype 5--7：
\begin{lstlisting}
-sel 1-5,0-7  6-8,5-7
\end{lstlisting}

排除边层 2，以及边层 4 的 datatype 5--6：
\begin{lstlisting}
-uel 2 4,5-6
\end{lstlisting}

% ============================================================
\section{文本处理}
\subsection*{Handling of Text}

\cmd{-text} 命令行选项指示 LVL 比较两个版图之间的文本。若一个或两个数据库在同一坐标 \cmd{(x1,y1)} 有多个文本，实用程序在 \cmd{(x1,y1)} 处尽可能匹配两库间的文本字符串。错误报告方式如下：

若每个数据库在 \cmd{(x1, y1)} 仅有一个未匹配文本：
\begin{lstlisting}
Violations Found:TEXT_XOR
- - - - - - - - - - - - - - - - - - - - - - - - -
Structure (Position x, y) BaselineText CheckText
- - - - - - - - - - - - - - - - - - - - - - - - -
A         (x1, y1)         base_text    check_text
\end{lstlisting}

若有一个未匹配基线文本且无未匹配检查文本（或相反）：
\begin{lstlisting}
Violations Found:TEXT_XOR
- - - - - - - - - - - - - - - - - - - - - - - - -
Structure (Position x, y) BaselineText CheckText
- - - - - - - - - - - - - - - - - - - - - - - - -
A         (x1, y1)         base_text
\end{lstlisting}

若基线数据库在 \cmd{(x1,y1)} 有多个未匹配文本，且检查库至少一个未匹配文本：
\begin{lstlisting}
Violations Found:TEXT_XOR
- - - - - - - - - - - - - - - - - - - - - - - - -
Structure (Position x, y) BaselineText CheckText
- - - - - - - - - - - - - - - - - - - - - - - - -
A         (x1, y1)         base_text1   check_text1
A         (x1, y1)         base_text1   check_text2
A         (x1, y1)         base_text2   check_text1
A         (x1, y1)         base_text2   check_text2
A         (x1, y1)         base_text3   check_text1
A         (x1, y1)         base_text3   check_text2
\end{lstlisting}

上例中有三个未匹配基线文本与两个未匹配检查文本。一般报告未匹配文本的全部可能组合。

% ============================================================
\section{LVL 输出}
\subsection*{LVL Output}

LVL 实用程序错误保存在「错误文件」所述文件中。

输出库格式与基线库相同。默认输出单元来自基线库。VUE 将差异叠加到基线库上。
\begin{itemize}
  \item 若输入为 GDSII 或 OASIS 且 LVL 使用层文件运行，差异报告在 \cmd{LVL\_ERRORS} 所列层上，可能与基线层或检查版图不同。例如，基线层 \cmd{\{31,0\}} 与检查层 \cmd{\{131,0\}}（经 LVL 选项映射到同一层）的差异可报告在输出层 \cmd{\{1,23\}}。
  \item 若输入为 GDSII 或 OASIS，使用层映射文件且无层文件，差异报告在由映射文件名称确定的层上。例如层名为 \cmd{L10D15} 时，差异报告在层 \cmd{\{10,15\}}。
  \item 若输入为 GDSII 或 OASIS，且既无层映射也无层文件，两层差异报告在输出库的同一层上。例如基线与检查均为 \cmd{\{31,0\}} 时，报告在输出层 \cmd{\{31,0\}}。
\end{itemize}

若某一输入文件（例如 OpenAccess 或 Milkyway）不是 GDSII 或 OASIS，差异报告在 \cmd{LVL\_ERRORS} 所列层上，可能与基线或检查版图不同。例如基线 \cmd{\{31,0\}} 与检查 \cmd{\{131,0\}}（经选项映射到同一层）的差异可报告在输出层 \cmd{\{1,23\}}。

\subsection{退出状态}
\subsubsection*{Exit Status}

\cmd{icv\_lvl} 可执行文件的退出状态见表~\ref{tab:app-b-exit}。

\begin{table}[htbp]
\centering
\caption{退出状态值}
\label{tab:app-b-exit}
\begin{tabular}{@{}lp{0.55\textwidth}@{}}
\toprule
退出状态值 & 定义 \\
\midrule
0 & 成功运行且未发现差异。 \\
1 & 成功运行且发现差异。 \\
负数 & 发生错误。 \\
\bottomrule
\end{tabular}
\end{table}

示例：成功且无差异：
\begin{lstlisting}
% icv_lvl lib1 lib2 -c cell
% echo $?
0
\end{lstlisting}

\subsection{错误文件}
\subsubsection*{Error Files}

表~\ref{tab:app-b-errfiles} 说明生成并存放在当前工作目录的错误文件。

\begin{noteBox}
可用 \cmd{-elpc} 指定每类检查的最大错误数。见表~\ref{tab:app-b-lvl-opts}。
\end{noteBox}

LVL 还提供屏幕输出，包含每次 XOR 或 NOT 运算的实时结果。

\begin{longtable}{@{}>{\ttfamily\raggedright\arraybackslash}p{0.38\textwidth} p{0.56\textwidth}@{}}
\caption{错误文件}\\
\label{tab:app-b-errfiles}\\
\toprule
\textbf{\rmfamily 文件名} & \textbf{说明} \\
\midrule
\endfirsthead
\toprule
\textbf{\rmfamily 文件名} & \textbf{说明} \\
\midrule
\endhead
\bottomrule
\endfoot

cell1\_vs\_cell2.LVL\_ERRORS &
按层分隔的全部比较错误的错误文件。 \\
\midrule

cell1\_vs\_cell2.LVL\_RESULTS &
结果文件，简要描述运行各步并指出是否有错误。有错误时，结果文件指向失败运行的 \cmd{run\_details} 目录中特定的 \cmd{cell.RESULTS} 文件。关于指定 \cmd{cell.RESULTS}，见 \textit{IC Validator Reference Manual} 中 \cmd{run\_options()}。 \\
\midrule

cell1\_vs\_cell2.gds\newline
cell1\_vs\_cell2.oasis\newline
cell1\_vs\_cell2/ &
含比较差异的输出库（GDSII、OASIS、Milkyway、OpenAccess、NDM 格式）。对 OpenAccess 库，会创建或更新当前目录的 \cmd{lib.defs} 以加入输出库定义。 \\
\midrule

lvl\_details/pydb &
SQLite 数据库（错误数据库 PYDB），按 LVL 执行的每次 XOR 或 NOT 运算分隔存放比较错误。 \\
\midrule

lvl\_details/cell1\_vs\_cell2.lvl\_sum &
摘要文件，含每次 LVL 运行的：\cmd{icv\_lvl} 命令行；成功或失败指示；运行目的；耗时与峰值内存。 \\
\midrule

cell1\_vs\_cell2.vue &
将差异叠加到基线库的文件，便于用 VUE 查看错误。 \\
\end{longtable}

\subsection{输出目录}
\subsubsection*{Output Directories}

表~\ref{tab:app-b-outdirs} 说明 LVL 在当前工作目录创建的目录。

\begin{longtable}{@{}>{\ttfamily\raggedright\arraybackslash}p{0.38\textwidth} p{0.56\textwidth}@{}}
\caption{输出目录}\\
\label{tab:app-b-outdirs}\\
\toprule
\textbf{\rmfamily 目录名} & \textbf{说明} \\
\midrule
\endfirsthead
\toprule
\textbf{\rmfamily 目录名} & \textbf{说明} \\
\midrule
\endhead
\bottomrule
\endfoot

lvl\_details/compare\_layers/ &
含基线库与检查库比较过程的额外诊断信息。 \\
\midrule

lvl\_details/extract\_layers/ &
含检查库预处理的摘要信息。 \\
\midrule

lvl\_details/layer\_list1/ &
若无层列表，为基线库自动生成的层列表；含层列表创建摘要供诊断。 \\
\midrule

lvl\_details/layer\_list2/ &
若无层列表，为检查库自动生成的层列表；含层列表创建摘要供诊断。 \\
\end{longtable}

% ============================================================
\section{使用模型}
\subsection*{Use Models}

本节给出 LVL 实用程序的两个简单使用模型。
\begin{itemize}
  \item 用两次 NOT 运算比较基线库与检查库；\cmd{library1} 的全部层与 \cmd{library2} 的全部层比较：
\begin{lstlisting}
icv_lvl library1 library2 -c cell -not
\end{lstlisting}
  \item 以 \cmd{cell} 为顶层单元，从输入库生成层文件；输出文件对单元子树中发现的每个 layer/datatype 含一条 assign 语句：
\begin{lstlisting}
icv_lvl library -af output_file -c cell
\end{lstlisting}
\end{itemize}

% ============================================================
\section{QuickLVL 实用程序}
\subsection*{QuickLVL Utility}

本节描述 Quick LVL 实用程序的两种使用模型，可快速读取两个版图并按单元逐一比较版图元素。
\begin{itemize}
  \item 单元级 Quick LVL：Quick LVL 的独立运行。
  \item 自动增量 DRC：IC Validator 先调用 Quick LVL 相对基线确定差异，结果以工具可读的内部格式传递，无需用户干预。
\end{itemize}

Quick LVL 支持下列版图输入格式：
\begin{itemize}
  \item OASIS
  \item GDSII
\end{itemize}

\begin{noteBox}
Quick LVL 实用程序支持除 NULL 字符外的任意文本字符。
\end{noteBox}

图~\ref{fig:app-b-quicklvl} 展示 Quick LVL 流程。

\figplaceholder{Figure 141: Quick LVL Flowchart}{Quick LVL 流程图}{fig:app-b-quicklvl}

\subsection{单元级 Quick LVL}
\subsubsection*{Cell-Level Quick LVL}

Quick LVL 读取两个版图，按单元逐一比较版图元素。结果存入错误数据库 PYDB。该模式生成 \cmd{LVL\_ERRORS} 与 VUE 格式输出文件。

命令行语法为：

\begin{lstlisting}
icv_lvl baseline_layout check_layout -c top_cell [other options] -quick
\end{lstlisting}

表~\ref{tab:app-b-lvl-opts} 说明 QuickLVL 的命令行选项。

\subsection{LVL 与 QuickLVL 的差异}
\subsubsection*{Differences Between LVL and QuickLVL}

单元级 Quick LVL 给出 \cmd{baseline\_layout} 与 \cmd{check\_layout} 之间的差异列表。与传统全版图 LVL 的差异如下：
\begin{itemize}
  \item 结果按单元级报告。
  \item 报告元素级差异，而非几何差异。例如某一矩形在一个版图中更大时，报告两个完整矩形元素，而非其几何差。
  \item 若命令行定义了变换，工具在尝试匹配基线版图前，先变换检查版图顶层单元中的坐标。
  \item 若同一元素在两个版图中以不同格式表示（例如 path 对 rectangle，或 AREF 对一组 SREF），即使几何上相同也不匹配，报告为“缺失”与“新增”元素。
  \item QuickLVL 在内部尝试匹配版图间元素。有一定相似性的元素报告为基线与检查之间的“不同”；其他元素报告为检查版图中缺失，或基线中的新形状。
  \item 可能出现一些误报（false-positive），以换取更高速度。
  \item QuickLVL 仅在单主机上工作。
\end{itemize}

该工具从 \cmd{LVL\_ERRORS} 文件提供如下输出示例：
\begin{lstlisting}
Path Mismatch (5:0)
  Baseline shapes different in Check ............. 25 violations found.

Placement Mismatch
  Baseline shapes missing in Check ............... 20 violations found.
  New shapes in Check not in Baseline ............ 150 violations found.

Path Mismatch (10:0)
  New shapes in Check not in Baseline ............ 75 violations found.
\end{lstlisting}

\subsection{聚焦差异层与区域的自动增量 DRC}
\subsubsection*{Auto-Incremental DRC Focused on Differing Layers and Regions}

使用 \cmd{-idrc\_i} 命令行选项时，IC Validator 先运行单元级 Quick LVL，将当前版图与基线版图比较。这些差异用于让工具聚焦含差异的层与窗口，从而改善运行时间。窗口 ambit 可通过设置 \cmd{incremental\_options(window\_ambit = \ldots);} 调整。

在自动增量 DRC 流程中，不推荐使用 \cmd{incremental\_options()} 中除 \cmd{window\_ambit} 以外的选项。该流程要求基线版图与检查版图具有相同的顶层单元名。自动增量 DRC 的结果可能无法完全替代完整的签核（signoff）运行。

自动增量 DRC 流程支持下列格式的基线版图：
\begin{itemize}
  \item OASIS
  \item GDSII
\end{itemize}

图~\ref{fig:app-b-idrc} 展示自动增量 DRC 流程。

\figplaceholder{Figure 142: Auto-Incremental DRC Flow}{自动增量 DRC 流程}{fig:app-b-idrc}
